\section{Experimentos}
\label{sec:experimentos}

\begin{table*}[!t]
\centering\small
\setlength{\tabcolsep}{4pt}
\begin{tabular}{lrccccccc}
\toprule
 & & \multicolumn{4}{c}{En distribución} & \multicolumn{2}{c}{OOD} & \\
\cmidrule(lr){3-6}\cmidrule(lr){7-8}
Modelo & Parám. (M) & N1 & N2 & N3 & N4 & N5 & N6 & Plant. nuevas \\
\midrule
GPT $L{=}1$ & 0.97 & $87.6{\scriptstyle\pm0.6}$ & $35.3{\scriptstyle\pm0.3}$ & $8.9{\scriptstyle\pm0.6}$ & $6.7{\scriptstyle\pm0.6}$ & $5.1{\scriptstyle\pm0.4}$ & $4.1{\scriptstyle\pm0.3}$ & $16.1{\scriptstyle\pm0.8}$ \\
GPT $L{=}2$ & 1.76 & $96.8{\scriptstyle\pm0.5}$ & $73.1{\scriptstyle\pm3.5}$ & $26.1{\scriptstyle\pm3.4}$ & $16.9{\scriptstyle\pm1.7}$ & $12.5{\scriptstyle\pm0.9}$ & $8.9{\scriptstyle\pm0.4}$ & $22.4{\scriptstyle\pm2.4}$ \\
GPT $L{=}4$ & 3.34 & $96.9{\scriptstyle\pm0.2}$ & $78.7{\scriptstyle\pm1.8}$ & $48.8{\scriptstyle\pm4.1}$ & $34.9{\scriptstyle\pm4.1}$ & $23.9{\scriptstyle\pm3.0}$ & $14.5{\scriptstyle\pm1.8}$ & $25.5{\scriptstyle\pm1.9}$ \\
GPT $L{=}8$ & 6.49 & $97.0{\scriptstyle\pm0.2}$ & $\mathbf{80.9}{\scriptstyle\pm0.4}$ & $54.3{\scriptstyle\pm2.1}$ & $40.2{\scriptstyle\pm1.8}$ & $25.4{\scriptstyle\pm2.7}$ & $13.6{\scriptstyle\pm1.7}$ & $26.7{\scriptstyle\pm0.6}$ \\
\midrule
Looped $N{=}2$ & 0.97 & $96.9{\scriptstyle\pm0.4}$ & $71.9{\scriptstyle\pm0.4}$ & $25.6{\scriptstyle\pm1.2}$ & $16.5{\scriptstyle\pm0.2}$ & $12.2{\scriptstyle\pm0.4}$ & $9.0{\scriptstyle\pm0.4}$ & $21.1{\scriptstyle\pm1.7}$ \\
Looped $N{=}4$ & 0.97 & $\mathbf{97.3}{\scriptstyle\pm0.2}$ & $78.7{\scriptstyle\pm0.6}$ & $47.0{\scriptstyle\pm1.1}$ & $31.6{\scriptstyle\pm0.3}$ & $21.1{\scriptstyle\pm2.3}$ & $14.0{\scriptstyle\pm0.1}$ & $23.3{\scriptstyle\pm1.0}$ \\
Looped $N{=}8$ & 0.97 & $97.1{\scriptstyle\pm0.2}$ & $79.0{\scriptstyle\pm0.2}$ & $\mathbf{57.2}{\scriptstyle\pm1.1}$ & $\mathbf{45.8}{\scriptstyle\pm1.4}$ & $\mathbf{33.4}{\scriptstyle\pm2.2}$ & $\mathbf{21.2}{\scriptstyle\pm2.2}$ & $\mathbf{29.8}{\scriptstyle\pm4.2}$ \\
\midrule
GPT $L{=}4$ ($d{=}128$) & 0.88 & $96.5{\scriptstyle\pm0.5}$ & $65.5{\scriptstyle\pm7.4}$ & $35.0{\scriptstyle\pm10.7}$ & $24.3{\scriptstyle\pm8.7}$ & $16.5{\scriptstyle\pm5.1}$ & $11.3{\scriptstyle\pm3.0}$ & $22.6{\scriptstyle\pm2.5}$ \\
\bottomrule
\end{tabular}

\caption{Exactitud (\%) por nivel, media $\pm$ desviación típica sobre tres
semillas (niveles 1--4: test en distribución; 5 y 6: fuera de distribución;
\emph{Plant. nuevas}: redacciones no vistas). En negrita, la mejor media.}
\label{tab:main}
\end{table*}

\begin{figure*}[!t]
\centering
\figorph{loops_sweep.pdf}{\textwidth}{exactitud frente al número de iteraciones usado en test}
\caption{Exactitud de un Looped Transformer entrenado con $N$ iteraciones al
evaluarlo con $N'$ sin reentrenar (banda: mínimo y máximo entre semillas).}
\label{fig:sweep}
\end{figure*}

\begin{figure*}[!t]
\centering
\figorph{internal_main.pdf}{\textwidth}{logit lens, atención a los dígitos y dinámica de representación y atención}
\caption{Interior de Looped $N{=}8$ y GPT $L{=}8$ con \emph{teacher forcing}.
Arriba: fracción de problemas con la respuesta completa correcta al proyectar
cada estado con la cabeza de salida (\emph{logit lens}) y atención de la posición
de respuesta a los dígitos del enunciado, por nivel. Abajo: cambio relativo,
rango efectivo, entropía normalizada y divergencia JS entre pasos (niveles 1--4;
banda: mínimo y máximo entre semillas).}
\label{fig:interno}
\end{figure*}

\subsection{Configuraciones y comparaciones controladas}

Todos los modelos usan contexto de 512 caracteres, vocabulario de 97 símbolos y
$d{=}256$ con 8 cabezas, salvo un GPT de 4 capas con $d{=}128$ y 4 cabezas
(contraste C3). Entrenamos GPT con
$L\in\{1,2,4,8\}$ capas y Looped Transformers con $N\in\{2,4,8\}$ iteraciones,
con las semillas 0, 1 y 2 (24 entrenamientos de 42\,000 pasos).
La Tabla~\ref{tab:main} indica el número de parámetros de cada modelo. La matriz
permite tres contrastes: \textbf{(C1)}~\emph{igual cómputo}, Looped $N$ frente a GPT
$L{=}N$, con los mismos FLOPs por token y $N$ veces más parámetros el segundo;
\textbf{(C2)}~\emph{iguales parámetros}, Looped $N$ frente a GPT $L{=}1$, que no
comparten el cómputo; y \textbf{(C3)}~\emph{parámetros y profundidad parecidos},
Looped $N{=}4$ (0.97\,M de parámetros) frente a ese GPT de 4 capas con $d{=}128$
(0.88\,M), que tiene la misma profundidad y casi los mismos parámetros pero 3.2
veces menos FLOPs por token.

\subsection{Resultados}

En las comparaciones, la primera cifra es siempre la del Looped y la segunda la
del GPT.

\paragraph{Igual número de parámetros (C2).}
Con 0.97\,M de parámetros, el Looped Transformer llega a 25.6 y 16.5\,\% en los
niveles 3 y 4 con $N{=}2$, a 47.0 y 31.6\,\% con $N{=}4$ y a 57.2 y 45.8\,\% con
$N{=}8$ (Tabla~\ref{tab:main}), mientras que el GPT de una capa apenas resuelve
esos niveles (8.9 y 6.7\,\%, frente al 6.2 y 5.4\,\% de responder siempre la
respuesta más frecuente): iterar el bloque aporta una capacidad que no se
obtiene con más parámetros.

\paragraph{Igual cómputo (C1).}
Con $N{=}2$ y $N{=}4$ el Looped queda al nivel del GPT de igual profundidad o
algo por debajo, dentro de la desviación entre semillas (nivel 4: 16.5 frente a
16.9\,\% con $N{=}L{=}2$, y 31.6 frente a 34.9\,\% con $N{=}L{=}4$), con 1.8 y 3.4
veces menos parámetros. Con $N{=}8$ iguala o supera al GPT de 8 capas, que tiene
6.7 veces más (nivel 3: 57.2 frente a 54.3\,\%; nivel 4: 45.8 frente a
40.2\,\%). A igual profundidad el tiempo de entrenamiento es el mismo (2.4 frente
a 2.5~h con $D{=}8$): el bucle ahorra parámetros, no tiempo. En C3, el Looped
$N{=}4$ supera al GPT con $d{=}128$ en los niveles 3 y 4 (47.0 frente a 35.0\,\% y
31.6 frente a 24.3\,\%), con mucha menos varianza entre semillas.

\paragraph{Generalización a más operaciones.}
En los niveles 5 y 6 la ventaja del Looped $N{=}8$ es mayor: 33.4 y 21.2\,\%
frente a 25.4 y 13.6\,\% del GPT de 8 capas, que apenas mejora al de 4 capas
(23.9 y 14.5\,\%) aunque sí lo supera en los niveles 3 y 4. Una posible causa es
que, con como máximo 4 operaciones en el entrenamiento, nada obliga a las capas
adicionales del GPT a repetir un mismo paso: ganan capacidad para ajustar los
niveles 1 a 4, mientras que el bucle debe aprender un paso reutilizable. Los
problemas OOD no son más largos que los de entrenamiento (411 caracteres como
máximo en ambos), así que no se trata de posiciones no vistas.

\paragraph{Número de iteraciones necesarias (P2).}
La Fig.~\ref{fig:sweep} varía las iteraciones en test sin reentrenar. Con menos
de las entrenadas la exactitud cae con rapidez, y antes cuanto más difícil es el
nivel: con $N{=}8$, el nivel 1 alcanza 89.6\,\% con 5 iteraciones y 97.1\,\% con 7,
mientras que el nivel 4 sigue mejorando hasta la octava (37.1\,\% con 7 y
45.8\,\% con 8). Las iteraciones necesarias crecen con la dificultad. Con
más de las entrenadas la exactitud también baja, y más deprisa cuanto menor es
$N$ ($N{=}4$ con 8 iteraciones: nivel 1 de 97.3 a 77.7\,\%), y en los niveles 5 y 6 el máximo está en $N$ o en $N{+}1$, por
lo que no observamos extrapolación en profundidad.
